
\documentclass{article}
\usepackage{colortbl}
\usepackage{makecell}
\usepackage{multirow}
\usepackage{supertabular}

\begin{document}

\newcounter{utterance}

\centering \large Interaction Transcript for game `eqbench', experiment `eq\_bench', episode 1 with reflexion{-}Qwen3.5{-}9B.
\vspace{24pt}

{ \footnotesize  \setcounter{utterance}{1}
\setlength{\tabcolsep}{0pt}
\begin{supertabular}{c@{$\;$}|p{.15\linewidth}@{}p{.15\linewidth}p{.15\linewidth}p{.15\linewidth}p{.15\linewidth}p{.15\linewidth}}
    \# & \multicolumn{2}{c}{Player} && \multicolumn{2}{c}{Game Master} \\
    \hline

    \theutterance \stepcounter{utterance}  
    & & & \multicolumn{4}{p{0.6\linewidth}}{
        \cellcolor[rgb]{0.9,0.9,0.9}{
            \makecell[{{p{\linewidth}}}]{
                \texttt{\tiny{[P1$\langle$GM]}}
                \texttt{You have attempted this same task before and did not succeed. These are the plans you wrote after those attempts.} \\
\texttt{Trial \#1: New plan: I incorrectly prioritized "Incredulous" and "Defensive" by focusing too heavily on Robert's aggressive tone and Claudia's sharp retorts, missing the underlying emotional context of their shared history and the specific threat to her life's work. I failed to recognize that while she is defending her values, the core of her reaction to Robert dismissing her parents' legacy as mere "sentimentality" is a deep sense of being misunderstood and a protective, almost sympathetic stance towards the nature they fought to preserve, rather than just anger at him. Next time, I will re{-}evaluate the dialogue by identifying the specific trigger for her emotional state: the devaluation of her parents' love. I will assign a higher score to "Defensive" (as she is actively countering his arguments) but significantly lower "Incredulous" (she believes him, she just disagrees) and higher "Sympathetic" (towards the trees/nature) or keep it low if the prompt implies sympathy for Robert, but actually, the most accurate reading is that she is **Defensive** of her principles and **Intimidated** by his cold pragmatism, while **Sympathetic** is low. Wait, looking closer: She calls him out on bulldozing, calls his world "this sanctuary," and says money won't bring back trees. She is standing her ground. "Defensive" is high. "Intimidated" is low because she is fighting back. "Incredulous" is low because she understands his drive but rejects it. "Sympathetic" is low. The mistake was overestimating disbelief. I will adjust to: Defensive: 9, Intimidated: 2, Incredulous: 1, Sympathetic: 1. Actually, let's look at the options again. Is she sympathetic? No. Is she incredulous? No, she's arguing. Is she intimidated? Maybe slightly by his power ("My company's survival"), but she holds ground. Is she defensive? Yes, protecting her "testament." The previous error was likely misreading the intensity of the conflict. I will ensure the highest score goes to the emotion that directly corresponds to her active rebuttal of his accusations.} \\
\\ 
\texttt{Revised Plan:} \\
\texttt{1. Analyze Claudia's specific lines: "Not my world," "testament to our parents' love," "money won't bring back the trees." These are defenses of her identity and values.} \\
\texttt{2. Recognize Robert's tone: "bulldozing," "progress," "survival." This is an attack on her values.} \\
\texttt{3. Determine the primary emotion: She is defending her stance against his pragmatism. "Defensive" is the strongest fit.} \\
\texttt{4. Evaluate secondary emotions: She is not "Incredulous" (she knows he is practical); she is not "Sympathetic"} \\
\\ 
\texttt{Trial \#2: New plan: I failed because I assigned the highest score to "Defensive" (9) while underestimating "Intimidated" (3). Upon re{-}evaluating the dialogue, Robert's shift from debating values to explicitly stating "My company's survival" and framing her sanctuary as an obstacle to "progress" creates a direct threat to her existence and livelihood, not just an argument about principles. While she is defending her values, the underlying pressure of his corporate power and the ultimatum nature of his argument ("survival") induces a stronger sense of fear and intimidation than pure defensiveness. She is standing her ground, but the *intensity* of the emotion driving her reaction is the threat to her world. Therefore, "Intimidated" should be the primary emotion, followed by "Defensive" as her reaction to that threat. "Incredulous" and "Sympathetic" remain low as she understands his pragmatism but rejects it. I will adjust the scores to reflect that the threat to her "world" makes intimidation the dominant feeling.} \\
\\ 
\texttt{Revised Scores:} \\
\texttt{Intimidated: 9} \\
\texttt{Defensive: 7} \\
\texttt{Incredulous: 1} \\
\texttt{Sympathetic: 1} \\
\\ 
\\ 
\\ 
\texttt{Your task is to predict the likely emotional responses of a character in this dialogue:} \\
\\ 
\texttt{Robert: Claudia, you've always been the idealist. But let's be practical for once, shall we?} \\
\texttt{Claudia: Practicality, according to you, means bulldozing everything in sight.} \\
\texttt{Robert: It's called progress, Claudia. It's how the world works.} \\
\texttt{Claudia: Not my world, Robert.} \\
\texttt{Robert: Your world? You mean this...this sanctuary of yours?} \\
\texttt{Claudia: It's more than a sanctuary. It's a testament to our parents' love for nature.} \\
\texttt{Robert: Sentimentality won't pay the bills, Claudia.} \\
\texttt{Claudia: And money won't bring back the trees you've felled, Robert.} \\
\texttt{Robert: This isn't about trees, Claudia. It's about survival. My company's survival.} \\
\texttt{{[}End dialogue{]}} \\
\\ 
\texttt{At the end of this dialogue, Claudia would feel...} \\
\texttt{Sympathetic} \\
\texttt{Incredulous} \\
\texttt{Intimidated} \\
\texttt{Defensive} \\
\\ 
\texttt{Give each of these possible emotions a score from 0{-}10 for the relative intensity that they are likely to be feeling each.} \\
\\ 
\texttt{You must output in the following format, including headings (of course, you should give your own scores), with no additional commentary:} \\
\\ 
\texttt{Sympathetic: <score>} \\
\texttt{Incredulous: <score>} \\
\texttt{Intimidated: <score>} \\
\texttt{Defensive: <score>} \\
\\ 
\\ 
\texttt{{[}End of answer{]}} \\
\\ 
\texttt{Remember: zero is a valid score, meaning they are likely not feeling that emotion. You must score at least one emotion > 0.} \\
\\ 
\texttt{Your answer:} \\
            }
        }
    }
     \\ \\

    \theutterance \stepcounter{utterance}  
    & \multicolumn{4}{p{0.6\linewidth}}{
        \cellcolor[rgb]{0.561,0.737,0.561}{
            \makecell[{{p{\linewidth}}}]{
                \texttt{\tiny{[P1$\rangle$GM]}}
                \texttt{Sympathetic: 1} \\
\texttt{Incredulous: 1} \\
\texttt{Intimidated: 2} \\
\texttt{Defensive: 9} \\
            }
        }
    }
    & & \\ \\

    \theutterance \stepcounter{utterance}  
    & & & \multicolumn{2}{p{0.3\linewidth}}{
        \cellcolor[rgb]{0.9,0.9,0.9}{
            \makecell[{{p{\linewidth}}}]{
                \texttt{\tiny{[GM$|$GM]}}
                \texttt{Sympathetic: 1} \\
\texttt{Incredulous: 1} \\
\texttt{Intimidated: 2} \\
\texttt{Defensive: 9} \\
            }
        }
    }
    & & \\ \\

    \theutterance \stepcounter{utterance}  
    & & & \multicolumn{2}{p{0.3\linewidth}}{
        \cellcolor[rgb]{0.9,0.9,0.9}{
            \makecell[{{p{\linewidth}}}]{
                \texttt{\tiny{[GM$|$GM]}}
                \texttt{Sympathetic: 0} \\
\texttt{Incredulous: 8} \\
\texttt{Intimidated: 3} \\
\texttt{Defensive: 7} \\
            }
        }
    }
    & & \\ \\

    \theutterance \stepcounter{utterance}  
    & & & \multicolumn{2}{p{0.3\linewidth}}{
        \cellcolor[rgb]{0.9,0.9,0.9}{
            \makecell[{{p{\linewidth}}}]{
                \texttt{\tiny{[GM$|$GM]}}
                \texttt{game\_result = LOSE} \\
            }
        }
    }
    & & \\ \\

    \theutterance \stepcounter{utterance}  
    & & \multicolumn{4}{p{0.6\linewidth}}{
        \cellcolor[rgb]{0.95,0.95,0.95}{
            \makecell[{{p{\linewidth}}}]{
                \texttt{\tiny{[P1$\leftrightarrow$P]}}
                \texttt{Post{-}episode Reflexion (not sent to GM):} \\
\texttt{New plan: I failed because I misidentified the primary emotional driver. In my previous attempts, I oscillated between "Defensive" and "Intimidated," but I overlooked the specific nuance of Robert's final line: "It's about survival. My company's survival." This shifts the dynamic from a debate of values to a direct threat of destruction. Claudia's response ("money won't bring back the trees") is not just a defense of principles; it is a desperate, futile attempt to argue against an inevitable loss of her life's work. The core emotion is not anger (Defensive) or fear (Intimidated), but a profound sense of helplessness and disbelief that her values can be so casually dismissed in favor of corporate gain. She is **Incredulous** that her "testament" is being treated as expendable. Therefore, I must assign the highest score to "Incredulous" to reflect her shock at the devaluation of her parents' legacy, while keeping "Defensive" moderate (she is arguing, but her argument is losing ground) and "Intimidated" low (she is fighting back, not cowering). "Sympathetic" remains low as she is focused on her own loss, not Robert's plight.} \\
\\ 
\texttt{Revised Scores:} \\
\texttt{Incredulous: 9} \\
\texttt{Defensive: 5} \\
\texttt{Intimidated: 2} \\
\texttt{Sympathetic: 1} \\
            }
        }
    }
    & \\ \\

\end{supertabular}
}

\end{document}
